\label{cap:fechamento}

Este trabalho investigou o problema de alocação de médicos em uma rede de clínicas de exames com múltiplas unidades autônomas, com base em dados reais de 250 profissionais e 37 unidades coletados ao longo de 2022 e 2023. O problema foi formulado como otimização inteira mista bi-objetiva, com minimização simultânea de exames não atendidos e trocas de clínica em horizontes acíclicos de planejamento, e resolvido por três famílias metodológicas distintas: o método exato por $\varepsilon$-restrito, MIP-heurísticas baseadas em \textit{Relax-and-Fix} e \textit{Fix-and-Optimize}, e meta-heurísticas evolutivas multiobjetivo nas variantes NSGA-II Estendido e Compacto.

As contribuições centrais do estudo podem ser sintetizadas em cinco pontos. Primeiro, a modelagem matemática inteira mista, de caráter bi-objetivo, para escalonamento médico em múltiplas clínicas autônomas, com tratamento conjunto da cobertura operacional e da estabilidade das alocações ao longo do horizonte. Segundo, a adaptação da variável de estado $\gamma_{iud}$ a partir do \textit{General Lot-sizing and Scheduling Problem} \cite{Fleischmann_Meyr_1997}, que viabiliza a detecção correta de trocas mesmo em agendas com lacunas. Terceiro, o acoplamento entre o método $\varepsilon$-restrito e mat-heurísticas para geração estruturada da fronteira de Pareto sob orçamento computacional comum. Quarto, a proposição das variantes NSGA-II Estendido e Compacto, equipadas com operadores customizados de amostragem com consciência de restrições, cruzamento por blocos diários e mutação por realocação de \textit{slot}. Por fim, a avaliação comparativa sistemática das três famílias metodológicas sobre dados reais, organizada em protocolo de dois estágios (parametrização e validação) ao longo de quatro horizontes de planejamento operacionalmente significativos.

\section{Principais Resultados}

A caracterização do método exato, apresentada no Capítulo~\ref{cap:resultados}, evidenciou a transição de regime entre os horizontes de 15 e 30 dias. No ciclo quinzenal, o Gurobi atinge otimalidade comprovada em parte significativa dos pontos $\varepsilon$, com gap MIP médio de $4{,}1\%$, e tempo total inferior ao limite de tempo alocado. A partir de 30 dias, o crescimento combinatório das variáveis $x_{iudt}$ inviabiliza a convergência: o gap médio sobe para $42{,}7\%$ aos 30 dias, $78{,}1\%$ aos 60 dias e $96{,}8\%$ aos 90 dias, com saturação completa do limite de 3.600 segundos por ponto $\varepsilon$. Esse padrão fundamenta a necessidade de heurísticas em horizontes operacionais maiores.

Entre as MIP-heurísticas, o método híbrido (Relax-and-Fix seguido de Fix-and-Optimize) apresentou desempenho próximo ao do método exato nos horizontes em que este converge, com hipervolume relativo de $0{,}996$ aos 15 dias e $1{,}008$ aos 30 dias. O Relax-and-Fix puro, embora ligeiramente inferior em qualidade ($0{,}984$ e $0{,}852$ nos mesmos horizontes), reduziu o tempo total de execução em até $30\%$ em relação ao exato, consolidando-se como alternativa viável em cenários com restrição temporal acentuada.

Nas meta-heurísticas, a representação compacta mostrou-se decisiva. O NSGA-II Compacto, ao codificar diretamente a associação diária médico-clínica, executa um número substancialmente maior de gerações dentro do limite de tempo, atingindo hipervolume relativo de $0{,}715$ contra $0{,}510$ da variante Estendida na fase de parametrização; no horizonte de 60 dias, o NSGA-IIc obteve HV relativo de $1{,}039$ em relação à referência exata --- que, neste horizonte, opera com MIP gap médio de $78{,}1\%$, tornando a comparação direta não conclusiva quanto à qualidade absoluta. A combinação da representação reduzida com os mecanismos de diversidade adaptativa, evolução em duas fases e inicialização enviesada para extremos, herdados da variante Estendida, configurou a abordagem evolutiva mais robusta entre as avaliadas.

Em conjunto, os resultados sustentam uma recomendação operacional clara para a rede de clínicas: o ciclo quinzenal de reprogramação pode ser tratado pelo método exato; horizontes mensal, bimestral e trimestral exigem necessariamente abordagens aproximadas, com o método híbrido como melhor escolha quando há tempo computacional comparável ao do \textit{solver} exato, e o NSGA-II Compacto como melhor escolha quando o limite de tempo é estritamente limitado.

\section{Limitações}

A principal limitação identificada reside no mecanismo de reparo de restrições do NSGA-II. A verificação estrita das soluções revelou que todas violam a restrição de cobertura mínima (\ref{eq:r5}) em 4 a 6\% das combinações clínica-período, inviabilizando o uso operacional direto sem inspeção e reparo posterior. A análise por variante revela distinções relevantes: no NSGA-II Estendido, as restrições estruturais \ref{eq:r1}--\ref{eq:r4} são integralmente satisfeitas e os valores de $Z_1$ e $Z_2$ foram verificados por Gurobi com variáveis fixadas sem discrepância; no NSGA-II Compacto, 10{,}2\% das soluções apresentam adicionalmente violação da restrição de capacidade de sala (\ref{eq:r4}), concentradas na clínica U06 no turno vespertino, e os valores de $Z_1$ reportados subestimam os exames não atendidos em até 23\% nas soluções no extremo de mínima cobertura --- comprometendo a validade dos objetivos para essa fração das soluções. Em consequência, nenhuma das variantes é diretamente comparável aos métodos MIP que garantem factibilidade, e os resultados do NSGA-II devem ser lidos isoladamente, dentro da família de meta-heurísticas.

A modelagem assume disponibilidade médica e demanda por exames como parâmetros determinísticos, obtidos a partir de registros históricos da rede de clínicas. Na operação real, ambas as grandezas estão sujeitas a variabilidade originada por absenteísmo, picos sazonais de demanda e reagendamentos, que o modelo atual absorve apenas indiretamente mediante otimização periódica em horizontes mais curtos. A variável de estado $\gamma_{iud}$ também exige uma condição inicial não nula para todos os médicos, o que, na ausência de histórico imediato, depende de inicialização aleatória ou de estimativa fornecida pelo gestor.

Permanecem fora da formulação algumas restrições ergonômicas presentes na operação cotidiana da rede: limites superiores ao número de finais de semana consecutivos trabalhados, número mínimo de domingos atribuídos por médico e descansos mandatórios entre turnos. Adicionalmente, a possibilidade de exceções negociadas às restrições~\eqref{eq:r2}-\eqref{eq:r3}, que atualmente impedem o profissional de atender múltiplas clínicas no mesmo dia, não está representada. Essas regras refletem acordos contratuais e operacionais que não foram incorporados para preservar a generalidade do modelo, mas constituem extensões naturais para implantação real.

\section{Trabalhos Futuros}

Dado o escopo abrangido neste trabalho, aliada a diversidade das possibilidades de extensão em pesquisas de Otimização e a riqueza nos resultados obtidos, torna-se natural que certas oportunidades não sejam abordadas. Nesse sentido, existe uma gama de direções investigativas em aberto que podem ser exploradas para aprofundar e ampliar os resultados aqui apresentados, tanto em termos de modelagem quanto de solução.

A princípio, pensa-se no desenvolvimento de operadores de reparo específicos para as limitações identificadas nas variantes do NSGA-II. Para a restrição de cobertura mínima (\ref{eq:r5}), violada em ambas as variantes, estratégias plausíveis incluem reparo \textit{post hoc} via subproblema MIP local, busca local guiada por unidades subcobertas e reformulação dos operadores de cruzamento e mutação para preservar a satisfação de \ref{eq:r5} por construção. Para a violação de capacidade de sala (\ref{eq:r4}) no NSGA-II Compacto, concentrada na clínica U06 no turno vespertino, a correção consiste em revisar o decodificador de cromossomo compacto para garantir o respeito aos limites de salas em todas as configurações de clínica e turno. A literatura recente em escalonamento na saúde \cite{Turhan_Bilgen_2020,Wickert_Kummer_Neto_Boniatti_Buriol_2021} indica que reparos baseados em MIPs locais combinam-se naturalmente com meta-heurísticas evolutivas, padrão coerente com a arquitetura adotada neste estudo.

Outra linha de investigação refere-se ao tratamento estocástico da demanda e da disponibilidade. Modelos robustos ou de programação estocástica em dois estágios podem absorver explicitamente a variabilidade observada na operação, complementados por mecanismos de reescalonamento \textit{online} que ajustem a escala diante de faltas de profissionais, picos de demanda e reagendamentos. A aplicação de técnicas de Aprendizado de Máquina apresenta-se como outra direção promissora em duas frentes complementares: a previsão da demanda por unidade-período, com vistas à substituição das estimativas históricas por projeções calibradas; e a geração de soluções iniciais informadas por padrões aprendidos, fornecendo \textit{warm-starts} de qualidade tanto para o método exato quanto para as mat-heurísticas e variantes do NSGA-II.

Por fim, a incorporação de uma metologia hibrida, aplicando conceitos de solução exata aliados a meta-heurísticas evolutivas em uma única \textit{pipeline} geracional pode explorar os limites computacionais de ambas as abordagens, reduzindo assim o \textit{gap} observado ao aplicar apenas uma metodologia individualizada. A combinação de métodos exatos e meta-heurísticas pode superar as limitações de cada abordagem isoladamente, especialmente em problemas de grande escala e alta complexidade como o escalonamento médico em múltiplas clínicas autônomas.